%\documentclass[12pt,letterpaper]{article}



%%%
%% Required packages:
%%
\usepackage{etex}
\usepackage{amsmath}
\usepackage{amssymb}
\usepackage{amstext}
\usepackage{amsthm}
\usepackage{fancyhdr,fancybox}
\usepackage{mathrsfs} % need for \mathscr command
\usepackage{multicol}
\usepackage[normalem]{ulem}
%\usepackage{pictex,pictexplus}
% \usepackage{pictexwd}
\usepackage{rotating}
\usepackage{url}
\usepackage{xfrac}
\usepackage{booktabs}
\usepackage{color,soul}
\usepackage{comment}

\usepackage{hyperref}
\hypersetup{
    colorlinks=true,     % false: boxed links; true: colored links
    linkcolor=blue,      % color of internal links
    citecolor=blue,      % color of links to bibliography
    filecolor=blue,      % color of file links
    urlcolor=blue        % color of external links
}


\usepackage{tikz}
\usetikzlibrary{backgrounds,calc}

%%
%% Common formatting commands
%%
\setlength{\parindent}{0in}
\setlength{\textwidth}{7in}
\setlength{\evensidemargin}{-0.25in}
\setlength{\oddsidemargin}{-0.25in}
\setlength{\parskip}{.5\baselineskip}
\setlength{\topmargin}{-0.5in}
\setlength{\textheight}{9in}

%               Problem and Part
\newcounter{problemnumber}
\newcounter{partnumber}
\newcounter{subpartnumber}
\newcommand{\Problem}{%
  \refstepcounter{problemnumber}%
  \setcounter{partnumber}{0}%
  \setcounter{subpartnumber}{0}%
  \item[\makebox{\hfil\textbf{\theproblemnumber.}\hfil}]%
  }
\newcommand{\InProblem}[2][1.6]{%
  \refstepcounter{problemnumber}%
  \setcounter{partnumber}{0}%
  \setcounter{subpartnumber}{0}%
  \textbf{\theproblemnumber.}\ \ \parbox[t]{#1in}{#2}%
}
\newcommand{\InSmallProblem}[1]{\InProblem[1.05]{#1}}
\newcommand{\NoProblem}[1]{%
  \mbox{\hfil\textbf{#1}\hfil}%
  }
\newcommand{\UnProblem}{%
  \refstepcounter{problemnumber}%
  \setcounter{partnumber}{0}%
  \setcounter{subpartnumber}{0}%
  \mbox{\hfil\textbf{\theproblemnumber.}\hfil}%
  }
\newcommand{\InFlexProblem}[1]{%
  \refstepcounter{problemnumber}%
  \setcounter{partnumber}{0}%
  \setcounter{subpartnumber}{0}%
  \textbf{\theproblemnumber.}\ \ #1%
}
%%  Part:
\newcommand{\Part}{%
  \renewcommand{\thepartnumber}{(\alph{partnumber})}%
  \refstepcounter{partnumber}
  \setcounter{subpartnumber}{0}%
  \item[(\alph{partnumber})]%
}
\newcommand{\InPart}[2][1.75]{%
  \renewcommand{\thepartnumber}{(\alph{partnumber})}%
  \refstepcounter{partnumber}%
  \setcounter{subpartnumber}{0}%
  (\alph{partnumber})\ \ \parbox[t]{#1in}{#2}%
}
\newcommand{\UnPart}{%
  \refstepcounter{partnumber}%
  \renewcommand{\thepartnumber}{(\alph{partnumber})}%
  \setcounter{subpartnumber}{0}%
  (\alph{partnumber})%
}
\newcommand{\InLargePart}[1]{\InPart[3.05]{#1}}
\newcommand{\InFullPart}[1]{\InPart[6.2]{#1}}
\newcommand{\InSmallPart}[1]{\InPart[1.05]{#1}}
%% SubPart:
\newcommand{\SubPart}{%
  \refstepcounter{subpartnumber}%
  \item[(\roman{subpartnumber})]%
  }
\newcommand{\InSubPart}[2][2.25]{%
  \stepcounter{subpartnumber}%
  (\roman{subpartnumber})\ \ \parbox[t]{#1in}{#2}%
}
\newcommand{\InSmallSubPart}[1]{\InSubPart[1.5]{#1}}
\newcommand{\InTinySubPart}[1]{%
  \stepcounter{subpartnumber}%
  (\roman{subpartnumber})\ \ \makebox{#1}%
}
\newcommand{\InMySubPart}[2][2.25]{%
  \stepcounter{subpartnumber}%
  \parbox[t]{#1in}{\makebox[0.3in][l]{(\roman{subpartnumber})}\ \ {#2}}%
}
\newcommand{\TrueFalse}{\textbf{T}\ \ \textbf{F}\ \ }
\newcommand{\TFPart}[1]{% 
  \stepcounter{partnumber}%
  \setcounter{subpartnumber}{0}%
  \item[(\alph{partnumber})] \TrueFalse\parbox[t]{5.5in}{#1}}

\newcommand{\Points}[1]{(#1 points) \ }
\newcommand{\Point}[1]{(#1 point) \ }


%% For Answers:
\newcounter{answernumber}
\newcommand{\Answer}{%
  \refstepcounter{answernumber}%
  \setcounter{partnumber}{0}%
  \setcounter{subpartnumber}{0}%
  \item[\makebox{\hfil\textbf{\theanswernumber.}\hfil}]%
  }


% Setting up the headers for the first page
%\chead{} % will default to what is typed in the file.
\fancyhead[L]{\textsc{Math 22a}}
\fancyhead[R]{\textsc{Fall 2024}}
\fancyfoot[R]{
	\vspace*{\fill}\footnotesize\textcopyright{} \emph{Harvard Math 22a}	
}
\renewcommand{\headrulewidth}{0pt}

%\fancyhead[C]{}
%	\chead{}	
%	\lhead{}
%	\rhead{}
%	\cfoot{}


% Putting copyright on the footer of each page
%\fancypagestyle{secondpage}{
%	\fancyhead[C]{TESTing again} % change this to be blank
%		%\chead{}	
%	\fancyhead[L]{bears} % change this to be blank
%	\fancyhead[R]{dogs} % change this to be blank
%	\fancyfoot[R]{ %keeps the footer the same
%		\vspace*{\fill}\footnotesize\textcopyright{} \emph{Harvard Math 22a}	
%	}
%}
%\renewcommand{\headrulewidth}{0pt}
%\pagestyle{fancy}
\pagestyle{empty}


%\lhead{}
%\rhead{}
%\rfoot{\vspace*{\fill}\footnotesize\textcopyright{} \emph{Harvard Math 22a}}
%\renewcommand{\headrulewidth}{0pt}
%\pagestyle{fancy}




%%
%% Common shortcut commands:
%%
\newcommand{\ds}{\displaystyle}
\newcommand{\sgn}{\operatorname{sgn}}
\renewcommand{\Pr}{\operatorname{P}}
\newcommand{\CPr}[3][{}]{\Pr_{\text{#1}}\left(#2\,|\,#3\right)}

\newcommand{\C}{\mathbb{C}}
\newcommand{\R}{\mathbb{R}}
\newcommand{\Z}{\mathbb{Z}}
\newcommand{\N}{\mathbb{N}}
\newcommand{\Q}{\mathbb{Q}}
\newcommand{\D}{\mathbb{D}}

\newcommand{\rref}{\operatorname{rref}}
\newcommand{\im}{\operatorname{im}}
\newcommand{\rank}{\operatorname{rank}}
\newcommand{\diag}{\operatorname{diag}}
\newcommand{\Span}{\operatorname{span}}
%\renewcommand{\vec}[1]{\mathbf{#1}}
\newcommand{\charpoly}{\mathscr{P}}
\newcommand{\nicefrac}[2]{\sfrac{#1}{#2}} % using xfrac package
\newcommand{\topology}{\mathcal{T}}
\newcommand{\Inn}{\operatorname{Inn}}

\newcommand{\blankbox}[2]{\fbox{\rule{#1}{0in}\rule{0in}{#2}}}

\newenvironment{myindentpar}[1]%
 {\begin{list}{}%
         {\setlength{\leftmargin}{#1}
          \setlength{\rightmargin}{\leftmargin}}%
         \item[]%
 }
 {\end{list}}
\newtheorem{theorem}{Theorem}
\newtheorem*{NStheorem}{Theorem}
\newtheorem{cor}[theorem]{Corollary}
\newtheorem*{claim}{Claim}
\newtheorem{defn}{Definition}

\newenvironment{amatrix}[1]{%
  \left(\begin{array}{@{}*{#1}{c}|c@{}}
}{%
  \end{array}\right)
}

%--------grstep
% For denoting a Gauss' reduction step.
% Use as: \grstep{\rho_1+\rho_3} or \grstep[2\rho_5 \\ 3\rho_6]{\rho_1+\rho_3}
\newcommand{\grstep}[2][\relax]{%
   \ensuremath{\mathrel{
       {\mathop{\longrightarrow}\limits^{#2\mathstrut}_{
                                     \begin{subarray}{l} #1 \end{subarray}}}}}}
\newcommand{\swap}{\leftrightarrow}



%\makeatletter
%\renewcommand*\env@matrix[1][*\c@MaxMatrixCols c]{%
%	\hskip -\arraycolsep
%	\let\@ifnextchar\new@ifnextchar
%	\array{#1}}
%\makeatother
%
%\def\env@matrix{\hskip -\arraycolsep
%	\let\@ifnextchar\new@ifnextchar
%	\array{*\c@MaxMatrixCols c}}
%

\section{{Determinants, $\vec\S$2.3}}% 
\chead{\Large\textbf{Determinants, $\vec\S$2.3}}% 

%\begin{document}
\thispagestyle{fancy}

Recall from the last few classes:

\begin{itemize}
\item We define the {\bf determinant} of $A=\begin{bmatrix} a & b \\ c & d\\ \end{bmatrix}$, denoted by $\text{det}(A)$, by $$\text{det}(A)=ad-bc.$$
\item {\bf Theorem:} If $A$ is an $n$ by $n$ invertible matrix, then for each $\vec{b}\in \R^n$, $A\vec{x}=\vec{b}$ has a unique solution; namely $\vec{x}=A^{-1}\vec{b}$.
\item {\bf Theorem:} Suppose $A$ and $B$ are $n$ by $n$ invertible matrices, then the following hold. 
\begin{enumerate}
\item $A^{-1}$ is invertible and $(A^{-1})^{-1}=A$.
\item $(AB)^{-1}=B^{-1} A^{-1}$.
\item $(A^{T})^{-1}=(A^{-1})^{T}$.
\end{enumerate}
\item An {\bf elementary matrix} $E$ is a matrix obtained from the identity matrix by a single elementary row operation.
 \item {\bf Theorem:} $A$ is invertible if and only if $A$ is row equivalent to the identity matrix. Furthermore, any sequence of elementary row operations that reduces $A$ to $I_n$ also transforms $I_n$ into $A^{-1}$.
\end{itemize}

\newpage





\begin{enumerate}

\Problem 
\Part Let $A=\begin{bmatrix}
1 & -2 & -1\\
-1 &5 & 6\\
5 &-4 & 5\\
\end{bmatrix}$. Find $A^{-1}$ if possible.

\vfill

\Part Let $A=\begin{bmatrix}
1 & 0 & -2\\
-3 &1 & 4\\
2 &-3 & 4\\
\end{bmatrix}$. Find $A^{-1}$ if possible.

\vfill

\newpage

{\bf Invertible Matrix Theorem} Let $A$ be an $n \times n$ matrix, then the following are all equivalent.
\begin{enumerate}
\item $A$ is invertible.
\item $A$ is row-equivalent to $I_n$.
\item $A$ has $n$ pivot positions.
\item $A\vec{x}=\vec{0}$ has only the trivial solution.
\item The columns of $A$ are linearly independent.
\item $x\mapsto A\vec{x}$ is one-to-one (injective).
\item For all $\vec{b}\in \R^n$, $A\vec{x}=\vec{b}$ is consistent.
\item The columns of $A$ span $\R^n$.
\item $\vec{x} \mapsto A\vec{x}$ is onto (surjective).
\item There exists a matrix $C$ such that $CA=I_n$.
\item There exists a matrix $D$ such that $AD=I_n$.
\item $A^{T}$ is invertible.
\end{enumerate}

\Problem Prove the theorem.

\newpage

\Problem Determine if the following matrices are invertible.
\Part Let $A=\begin{bmatrix}
3 & 0 & 0\\
-3 &-4 & 0\\
8 & 5 & -3\\
\end{bmatrix}$.

\vfill

\Part Let $A=\begin{bmatrix}
1 &2 & 3\\
4 &5 & 6\\
7 & 8 & 9\\
\end{bmatrix}$.

\vfill

\newpage

{\bf Theorem:} Let $T: \R^n \to \R^n$ be linear and let $A$ be the associated matrix. Then $T$ is invertible if and only if $A$ is invertible.

\Problem Prove the theorem.

\newpage

{\bf Goal:} Develop a numerical test for invertibility.

\Problem Let $A=\begin{bmatrix}
a & b & c\\
d & e &f\\
g & h & i\\
\end{bmatrix}$. Find a condition on the entries of $A$ for $A$ to be invertible.

\newpage

\begin{defn}
For a matrix $A$ define  $A_{ij}$ to be the matrix obtained from $A$ by deleting row $i$ and column $j$.
\end{defn}

{\bf Remark:} $$\text{det} \begin{bmatrix}
a & b & c\\
d & e &f\\
g & h & i\\
\end{bmatrix} = a \text{det} A_{11} - b \text{det} A_{12} +c \text{det} A_{13}.$$

\begin{defn}
For $n\geq 2$, the determinant of an $n\times n$ matrix $A=[a_{ij}]$ is
\begin{equation*}
\operatorname{det}(A)= a_{11} \operatorname{det}(A_{11}) - a_{12} \operatorname{det}(A_{12})+\dots + (-1)^{n+1} a_{1n} \operatorname{det}(A_{1n})=\sum_{j=1}^{n}(-1)^{1+j}a_{1j} \operatorname{det}(A_{1j}).
\end{equation*}
\end{defn}

\Problem Compute the determinant of $A=\begin{bmatrix}
1 & 2 &3\\
4 & 5 &6\\
7 & 8 & 9\\
\end{bmatrix}$.

\vfill

\begin{defn} The {\bf $(i,j)$-cofactor} of $A$ is given by $c_{ij} =(-1)^{i+j} \operatorname{det} A_{ij}$, and the {\bf cofactor expansion along row $i$} is given by $\sum_{j=1}^{n}a_{ij} c_{ij}$.
\end{defn}



{\bf Theorem 3.1:} Cofactor expansion across any row or column is the same, and hence each cofactor expansion gives the determinant.

\newpage

\Problem Compute the determinant of 
$C=\begin{bmatrix}
6 & 0 & 0 & 5\\
1 & 7 & 2 & -5\\
2 & 0 & 0 & 0\\
8 & 3 & 1 & 8\\
\end{bmatrix}$.

\vfill

\begin{defn} Let $A=[a_{ij}]$ be an $n\times n$ matrix, then $A$ is {\bf upper triangular} if $a_{ij}=0$ when $i>j$.
\end{defn}

{\bf Theorem 3.2:} Let $A$ be an $n\times n$ upper triangular matrix, then $$\operatorname{det}(A)= a_{11}a_{22}\dots a_{nn} = \prod_{j=1}^n a_{jj}.$$

\Problem Prove the theorem.

\vfill
\vfill

\end{enumerate}

\begin{comment}
%\end{document}

\newpage
\chead{\Large\textbf{Determinants -- Answers / Solutions}}%
\lhead{}
\rhead{}
\thispagestyle{fancy}

\begin{enumerate}
  \Answer
  
 \begin{enumerate} 
  \Part Let $A=\begin{bmatrix}
  1 & -2 & -1\\
  -1 &5 & 6\\
  5 &-4 & 5\\
  \end{bmatrix}$. We try augment $A$ by $I_n$ and try to use the result that if $A$ is invertible, then $[A|I_n]\sim [I_n | A^{-1}]$. Thus
  
  \begin{align*}
  \begin{pmatrix}[ccc|ccc]
  1 & -2 & -1 & 1 & 0 & 0\\
  -1 &5 & 6 & 0 & 1 & 0\\
  5 &-4 & 5 & 0 & 0 &1\\
  \end{pmatrix}& \sim  \begin{pmatrix}[ccc|ccc]
  1 & -2 & -1 & 1 & 0 & 0\\
  0 &3 & 5 & 1 & 1 & 0\\
  0 & 6 & 10 & -5 & 0 &1\\
\end{pmatrix} \sim \begin{pmatrix}[ccc|ccc]
1 & -2 & -1 & 1 & 0 & 0\\
0 &3 & 5 & 1 & 1 & 0\\
0 &0 & 0 & -7 & -2 &1\\
\end{pmatrix},
  \end{align*}
and now we're stuck. Since $A$ is not row-equivalent to $I_n$ (which we're seeing on left), the algorithm fails to find an inverse for $A$. Thus we know $A$ is not invertible.

\Part Let  $A=\begin{bmatrix}
1 & 0 & -2\\
-3 &1 & 4\\
2 &-3 & 4\\
\end{bmatrix}$. We again augment $A$ by $I_n$ and try to use the result that if $A$ is invertible, then $[A|I_n]\sim [I_n | A^{-1}]$. Thus

\begin{align*}
	\begin{pmatrix}[ccc|ccc]
		1 & 0 & -2 & 1 & 0 & 0\\
		-3 &1 & 4 & 0 & 1 & 0\\
		2 &-3 & 4 & 0 & 0 &1\\
	\end{pmatrix}& \sim  \begin{pmatrix}[ccc|ccc]
		1 & 0 & -2 & 1 & 0 & 0\\
		0 &1 & -2 & 3 & 1 & 0\\
		0 & -3 & 8 & -2 & 0 &1\\
	\end{pmatrix}\\ &\sim \begin{pmatrix}[ccc|ccc]
		1 & 0 & -2 & 1 & 0 & 0\\
		0 &1 & -2 & 3 & 1 & 0\\
		0 &0 & 2 & 7 & 3 &1\\
	\end{pmatrix}\\ &\sim
\begin{pmatrix}[ccc|ccc]
	1 & 0 & 0 & 8 & 3 &1 \\
	0 &1 & 0 & 10 & 4 & 1\\
	0 &0 & 2 & \frac{7}{2} & \frac{3}{2} &\frac{1}{2}\\
\end{pmatrix},
\end{align*}
and hence $A^{-1}=\begin{bmatrix}
8 & 3 & 1\\
10 & 4 & 1\\
\frac{7}{2} & \frac{3}{2} & \frac{1}{2}\\
\end{bmatrix}.$

\end{enumerate}
  
  \Answer 
  
 \begin{proof} There are many ways to organize all the implications. Here I try to use previous results from class as much as possible.
 
 $(a)\implies (j)$ by definition.
 
 $(j) \implies (d)$ by Theorem 2.5 (or above on this handout).
 
 $(d) \implies (c)$ by Theorem 1.4.
 
 $(c) \implies (b)$ since $A$ is square.
 
 $(b) \implies (a)$ by Theorem 2.7 (or above on this handout).
 
 Thus we have established $(a) \iff (b) \iff (c) \iff (d) \iff (j)$.
 
Note $(a)\implies (k)$ by definition of invertibility. We prove $(k) \implies (g)$. Suppose $D$ is an $n\times n$ matrix such that $AD=I_n$. Let $\vec{b} \in \mathbb{R}^n$, and define $\vec{y} =D\vec{b}$. Now $A\vec{y}=AD\vec{b} =I_n\vec{b} =\vec{b}$, and hence $\vec{y}$ is a solution to the matrix equation $A\vec{x}=\vec{b}$. Thus the matrix equation is always consistent, and hence $(g)$ follows. Now we prove $(g)\implies (a)$. By Theorem 1.4, $A\vec{x}=\vec{b}$ being consistent for any $\vec{b}$ is equivalent to $A$ having a pivot position in each row. Since $A$ is square, we know $A$ must be row equivalent to the identity matrix. By the theorem above, $A$ is invertible, and hence $(a)$ holds. 

Finally note $(d) \iff (e) \iff (f)$ and $(g) \iff (h) \iff (i)$ follows from Theorem 1.12. Finally we prove $(a)\iff (l)$ above on this handout.
 \end{proof}
  
  \Answer 
  
  \begin{enumerate}
  \Part We could row-reduce $A$ to determine if $A$ is row equivalent to the identity matrix, but in this case we can see the columns of $A$ are linearly independent, and hence we know $A$ must be invertible from the invertible matrix theorem.
  
  \Part We have previously seen this matrix does not span all of $\mathbb{R}^3$, and hence we can use the invertible matrix theorem to conclude $A$ is not invertible.
  
  \end{enumerate}
  
  \Answer
  
  \begin{proof} We prove both implications directly from the definitions.
  
  $\implies$ Suppose $T$ is invertible, then there exists a function $T^{-1} : \mathbb{R}^n \to \mathbb{R}^n$ such that $T(T^{-1}(\vec{x}))=\vec{x}$ and $T^{-1}(T(\vec{x}))=\vec{x}$ for all $\vec{x} \in \mathbb{R}^n$. If we knew that $T^{-1}$ was linear, then it would have an associated matrix we could use.
  
  {\bf Claim:} $T^{-1}$ is linear.
  
  {\bf Proof of Claim:} Let $\vec{x}, \vec{y} \in \mathbb{R}^n$ and $c \in \mathbb{R}$. Since $T$ is invertible, it is bijective. Thus there exists unique $\vec{a}, \vec{b} \in \mathbb{R}^n$ such that $T(\vec{a})=\vec{x}$ and $T(\vec{b})=\vec{y}$. Now we check the conditions for linearity: $$T^{-1}(\vec{x}+\vec{y})=T^{-1}(T(\vec{a})+T(\vec{b}))=T^{-1}(T(\vec{a}+\vec{b})=\vec{a}+\vec{b}=T^{-1}(\vec{x})+T^{-1}(\vec{y}).$$ and $$T^{-1}(c\vec{x})=T^{-1}(c T(\vec{a}))=T^{-1}(T(c\vec{a}))=c \vec{a} = c T^{-1}(\vec{x}).$$ Thus $T^{-1}$ is linear.
  
  Since $T^{-1}$ is linear, it has an associated matrix $D$. Now $$AD\vec{x} =T(T^{-1}(\vec{x}))=\vec{x}$$ and $$DA\vec{x} =T^{-1}(T(\vec{x})) =\vec{x}.$$ Thus $D$ is the inverse matrix of $A$, and hence $A$ is invertible. 
  
  $\impliedby$ Suppose $A$ is invertible. Define $S: \mathbb{R}^n \to \mathbb{R}^n$ by $S(\vec{x})=A^{-1}\vec{x}$. Then $T(S(\vec{x}))=AA^{-1}\vec{x} = \vec{x}$ and $S(T(\vec{x})) =A^{-1}A\vec{x}=\vec{x}$. Thus $S$ is the inverse function of $T$, and hence $T$ is invertible. 
  \end{proof}


  \Answer By the invertible matrix theorem, we know $A$ is invertible if and only if $A$ is row equivalent to the identity matrix. Thus we try to row-reduce $A$ and see what condition we get on the entries of the matrix.
  
  If the first column of $A$ is all zeros, then we know the matrix is not invertible. Thus we assume $a \neq 0$.
  
  \begin{align*}
  A=\begin{bmatrix}
  a & b & c\\
  d & e & f\\
  g & h & i\\
  \end{bmatrix}&\sim
\begin{bmatrix}
	a & b & c\\
	ad & ae & af\\
	ag & ah & ai\\
\end{bmatrix}\sim
\begin{bmatrix}
	a & b & c\\
	0 & ae-bd & af-cd\\
	0 & ah-bg & ai-cg\\
\end{bmatrix}\\
&\sim
\begin{bmatrix}
	a & b & c\\
	0 & ae-bd & af-cd\\
	0 & (ah-bg)(ae-bd) & (ai-cg)(ae-bd)\\
\end{bmatrix}\\
&\sim\begin{bmatrix}
	a & b & c\\
	0 & ae-bd & af-cd\\
	0 & 0 & D\\
\end{bmatrix} 
  \end{align*}
where $D=(ai-cg)(ae-bd)-(af-cd)(ah-bg)$.

Now multiplying out $D$ and simplifying we get 
$$D=a(aei-ceg-bdi -afh+cdh+bfg).$$ Since $a$ is non-zero, we know for $A$ to be row equivalent to the identity matrix, we need $(aei-acg-bdi -afh+cdh+bfg)\neq0$. We define $$\det(A)=(aei-acg-bdi -afh+cdh+bfg).$$ 

This seems like a strange condition, so let's play around with it a bit more. Factoring we get the following $$(aei-acg-bdi -afh+cdh+bfg)= a(ei-fh)-b(di-fg)+c(dh-eg),$$ and now we see some smaller determinants showing up; namely $$a(ei-fh)-b(di-fg)+c(dh-eg)= a \cdot \det \begin{bmatrix} e & f\\ h & i \end{bmatrix} -b\cdot \det \begin{bmatrix} d & f\\ g & i \end{bmatrix}+c\cdot \det \begin{bmatrix} d & e\\ g & h \end{bmatrix}=\det(A).$$ This expression for the determinants suggests a way to define the determinant of an $n\times n$ matrix in terms of $(n-1)\times(n-1)$ matrices.
  
  \Answer We use the definition to compute the determinant:
  \begin{align*}
 \det A&=\det \begin{bmatrix}
  1 & 2 &3\\
  4 & 5 &6\\
  7 & 8 & 9\\
  \end{bmatrix}=1 \cdot \det \begin{bmatrix} 5 & 6\\ 8 & 9 \end{bmatrix}-2 \cdot \det \begin{bmatrix} 4 & 6\\ 7 & 9 \end{bmatrix}+3 \cdot \det \begin{bmatrix} 4 & 5\\ 7 & 8 \end{bmatrix}\\&=1(45-48)-2(36-42)+3(32-35)=0.
\end{align*}

We have previously seen that $A$ is not invertible, so it is reassuring to see the determinant is also zero.
    
  \Answer Again we use the definition to compute the determinant:
  \begin{align*}
  	\det C& =\det \begin{bmatrix}
  6 & 0 & 0 & 5\\
  1 & 7 & 2 & -5\\
  2 & 0 & 0 & 0\\
  8 & 3 & 1 & 8\\
  \end{bmatrix}= 6\cdot \det \begin{bmatrix}
  7 & 2 & -5\\
  0 & 0 & 0\\
  3 & 1 & 8\\
\end{bmatrix}-5\cdot \det \begin{bmatrix}
1 & 7 & 2\\
2 & 0 & 0\\
8 & 3 & 1\\
\end{bmatrix}=-5(-2)=10.
\end{align*}

Using the last Theorem, we can expand about the third row as well. We show that here:
  \begin{align*}
	\det C& =\det \begin{bmatrix}
		6 & 0 & 0 & 5\\
		1 & 7 & 2 & -5\\
		2 & 0 & 0 & 0\\
		8 & 3 & 1 & 8\\
	\end{bmatrix}= 2\cdot \det \begin{bmatrix}
		0 & 0 & 5\\
		7 & 2 & -5\\
		3 & 1 & 8\\
\end{bmatrix}
=2(5)(7-6)=10.
\end{align*}


  \Answer 
  
  \begin{proof}
  We prove the result by induction on the size of the matrix.
  
  {\bf Base Case:} Let $A=\begin{bmatrix} a & b\\ 0 & c\\ \end{bmatrix}$, then $\det(A)=ac$. Thus the result holds for $2\times 2$ matrices.
  
  {\bf Inductive Step:} We assume that the determinant of an $n\times n$ upper triangular matrix is given by the product of the diagonal elements.
  
  Let $A$ be an $(n+1)\times(n+1)$ upper triangular matrix. We compute the determinant by expanding along the first column:
  \begin{align*}
  \det(A)&= a_{11} \det(A_{11}) -a_{21}\det(A_{21})+\dots+(-1)^{n+1+1}a_{(n+1)1}\det(A_{(n+1)1})
  \end{align*}

Since $A$ is upper-triangular, $a_{21}=\dots=a_{(n+1)1}=0$. Thus
$$\det(A)= a_{11} \det(A_{11}).$$ Since $A_11$ is an $n\times n$ upper triangular matrix, we can apply the induction hypothesis to get $$\det(A)= a_{11} \det(A_{11})=a_{11} (a_{22}\dots a_{(n+1)(n+1)}).$$ Thus the result holds by mathematical induction.
  \end{proof}
  
\end{enumerate}

\end{comment}
%\end{document}


%%% Local ables: 
%%% mode: latex
%%% TeX-master: t
%%% End: 
